\documentclass[12pt,a4paper]{article}
\usepackage[utf8]{inputenc}
\usepackage[vietnamese]{babel}
\usepackage{amsmath,amssymb,amsthm}
\usepackage{geometry}
\usepackage{enumitem}
\usepackage[most]{tcolorbox}
\usepackage{hyperref}
\usepackage{fancyhdr}
\usepackage{tikz}
\usepackage{eso-pic}
\usepackage{xcolor}
\geometry{a4paper, margin=2cm, bottom=2.5cm}
\hypersetup{colorlinks=true, linkcolor=blue!70!black}
\setlist[itemize]{topsep=4pt, itemsep=2pt}
\setlist[enumerate]{topsep=4pt, itemsep=2pt}
\AddToShipoutPictureFG{%
  \AtPageCenter{%
    \begin{tikzpicture}[remember picture, overlay]
      \node[rotate=45, scale=1, opacity=0.06]{\fontsize{60}{72}\selectfont\bfseries\sffamily Đặng Tấn Quí};
    \end{tikzpicture}%
  }%
}
\pagestyle{fancy}
\fancyhf{}
\fancyhead[L]{\small\textit{GIẢI TÍCH HÀM}}
\fancyhead[R]{\small\textit{Đặng Tấn Quí}}
\fancyfoot[C]{\thepage}
\fancyfoot[L]{\footnotesize\textcopyright\ Đặng Tấn Quí}
\fancyfoot[R]{\footnotesize SĐT: 0377\,122\,210}
\renewcommand{\headrulewidth}{0.4pt}
\renewcommand{\footrulewidth}{0.4pt}
\fancypagestyle{plain}{\fancyhf{}\fancyfoot[C]{\thepage}\fancyfoot[L]{\footnotesize\textcopyright\ Đặng Tấn Quí}\fancyfoot[R]{\footnotesize SĐT: 0377\,122\,210}\renewcommand{\headrulewidth}{0pt}\renewcommand{\footrulewidth}{0.4pt}}
\newtcolorbox{dinhnghia}[1][]{enhanced, breakable, colback=blue!3!white, colframe=blue!60!black, fonttitle=\bfseries, title={Định nghĩa}, #1}
\newtcolorbox{dinhly}[1][]{enhanced, breakable, colback=green!3!white, colframe=green!50!black, fonttitle=\bfseries, title={Định lý}, #1}
\newtcolorbox{tinhchat}[1][]{enhanced, breakable, colback=yellow!5!white, colframe=orange!50!black, fonttitle=\bfseries, title={Tính chất}, #1}
\newtcolorbox{phuongphap}[1][]{enhanced, breakable, colback=gray!5!white, colframe=gray!50!black, fonttitle=\bfseries, title={Phương pháp}, #1}
\newtcolorbox{luuy}[1][]{enhanced, breakable, colback=red!3!white, colframe=red!50!black, fonttitle=\bfseries, title={Lưu ý}, #1}
\newtcolorbox{congthuc}[1][]{enhanced, breakable, colback=cyan!3!white, colframe=cyan!50!black, fonttitle=\bfseries, title={Công thức}, #1}
\newtcolorbox{vidu}[1][]{enhanced, breakable, colback=violet!3!white, colframe=violet!50!black, fonttitle=\bfseries, title={Ví dụ}, #1}

\begin{document}
\thispagestyle{empty}
\begin{tikzpicture}[remember picture, overlay]
  \fill[blue!80!black] (current page.south west) rectangle (current page.north east);
  \node[anchor=west, white, font=\fontsize{26}{32}\bfseries\selectfont]
    at ([xshift=3cm, yshift=-7.5cm]current page.north west) {GIẢI TÍCH HÀM};
  \node[anchor=west, white, font=\fontsize{18}{24}\selectfont]
    at ([xshift=3cm, yshift=-9cm]current page.north west) {Không gian Banach, Hilbert, toán tử, lý thuyết phổ};
  \node[anchor=west, white, font=\Large\bfseries] at ([xshift=3cm, yshift=-13cm]current page.north west) {Biên soạn:};
  \node[anchor=west, yellow!80!white, font=\fontsize{22}{28}\bfseries\selectfont] at ([xshift=3cm, yshift=-14.5cm]current page.north west) {ĐẶNG TẤN QUÍ};
  \node[anchor=south east, white, font=\Large\bfseries] at ([xshift=-2cm, yshift=3cm]current page.south east) {\the\year};
\end{tikzpicture}
\newpage
\tableofcontents
\newpage
%% ================================================================
%%  CHƯƠNG 1: KHÔNG GIAN METRIC
%% ================================================================
\section{Không gian metric}

\begin{dinhnghia}
  \textbf{Không gian metric} $(X, d)$: tập $X$ với hàm khoảng cách $d: X \times X \to [0, +\infty)$ thỏa:
  \begin{enumerate}
    \item $d(x, y) = 0 \Leftrightarrow x = y$.
    \item $d(x, y) = d(y, x)$ (đối xứng).
    \item $d(x, z) \le d(x, y) + d(y, z)$ (bất đẳng thức tam giác).
  \end{enumerate}
\end{dinhnghia}

\begin{dinhnghia}[title={Sự hội tụ và tính đầy đủ}]
  \begin{itemize}
    \item Dãy $(x_n) \to x$ nếu $d(x_n, x) \to 0$.
    \item \textbf{Dãy Cauchy:} $\forall \varepsilon > 0$, $\exists N$: $m, n > N \Rightarrow d(x_m, x_n) < \varepsilon$.
    \item $(X, d)$ \textbf{đầy đủ} nếu mọi dãy Cauchy đều hội tụ trong $X$.
  \end{itemize}
\end{dinhnghia}

\begin{dinhnghia}[title={Tập compact}]
  $K \subset X$ \textbf{compact} nếu mọi phủ mở đều có phủ con hữu hạn. Tương đương (trong metric): mọi dãy trong $K$ có dãy con hội tụ trong $K$ (\textbf{compact dãy}).
\end{dinhnghia}

\begin{dinhly}[title={Ánh xạ co --- Định lý Banach (điểm bất động)}]
  Cho $(X, d)$ đầy đủ và $T: X \to X$ là \textbf{ánh xạ co}: $d(Tx, Ty) \le q\,d(x, y)$ với $0 \le q < 1$. Thì $T$ có duy nhất điểm bất động $x^* = Tx^*$ và $x_n = T^n(x_0) \to x^*$ với mọi $x_0$.
\end{dinhly}

%% ================================================================
%%  CHƯƠNG 2: KHÔNG GIAN ĐỊNH CHUẨN VÀ BANACH
%% ================================================================
\section{Không gian định chuẩn và Banach}

\begin{dinhnghia}[title={Không gian định chuẩn}]
  $(X, \|\cdot\|)$: $X$ là không gian vectơ (trên $\mathbb{R}$ hoặc $\mathbb{C}$) với chuẩn $\|\cdot\|: X \to [0, +\infty)$ thỏa:
  \begin{enumerate}
    \item $\|x\| = 0 \Leftrightarrow x = 0$.
    \item $\|\alpha x\| = |\alpha|\,\|x\|$ (thuần nhất).
    \item $\|x + y\| \le \|x\| + \|y\|$ (bất đẳng thức tam giác).
  \end{enumerate}
  Chuẩn sinh metric: $d(x, y) = \|x - y\|$.
\end{dinhnghia}

\begin{dinhnghia}[title={Không gian Banach}]
  Không gian định chuẩn \textbf{đầy đủ} gọi là \textbf{không gian Banach}.
\end{dinhnghia}

\begin{vidu}
  \begin{itemize}
    \item $\mathbb{R}^n$ với $\|x\|_p = \left(\sum |x_i|^p\right)^{1/p}$ ($1 \le p < \infty$) hoặc $\|x\|_\infty = \max|x_i|$: Banach.
    \item $C[a, b]$ với $\|f\|_\infty = \max_{[a,b]}|f(x)|$: Banach.
    \item $L^p(\mu)$ với $\|f\|_p = \left(\int |f|^p\,d\mu\right)^{1/p}$: Banach ($1 \le p \le \infty$).
    \item $\ell^p = \{(x_n) : \sum |x_n|^p < \infty\}$: Banach.
  \end{itemize}
\end{vidu}

\begin{dinhnghia}[title={Chuỗi trong không gian Banach}]
  Chuỗi $\sum x_n$ trong $X$ \textbf{hội tụ tuyệt đối} nếu $\sum \|x_n\| < \infty$.
\end{dinhnghia}

\begin{dinhly}
  $X$ là Banach khi và chỉ khi mọi chuỗi hội tụ tuyệt đối đều hội tụ.
\end{dinhly}

\begin{dinhnghia}[title={Chuẩn tương đương}]
  $\|\cdot\|_1$ và $\|\cdot\|_2$ \textbf{tương đương} nếu $\exists c, C > 0$: $c\|x\|_1 \le \|x\|_2 \le C\|x\|_1$, $\forall x$.
\end{dinhnghia}

\begin{dinhly}
  Trên không gian hữu hạn chiều, mọi chuẩn đều tương đương. Mọi không gian định chuẩn hữu hạn chiều đều là Banach.
\end{dinhly}

%% ================================================================
%%  CHƯƠNG 3: TOÁN TỬ TUYẾN TÍNH BỊ CHẶN
%% ================================================================
\section{Toán tử tuyến tính bị chặn}

\begin{dinhnghia}
  $T: X \to Y$ ($X, Y$ là không gian định chuẩn) là \textbf{toán tử tuyến tính} nếu $T(\alpha x + \beta y) = \alpha Tx + \beta Ty$.

  $T$ \textbf{bị chặn} nếu $\exists C \ge 0$: $\|Tx\| \le C\|x\|$, $\forall x \in X$.
\end{dinhnghia}

\begin{dinhly}
  $T$ tuyến tính thì: $T$ bị chặn $\Leftrightarrow$ $T$ liên tục $\Leftrightarrow$ $T$ liên tục tại $0$.
\end{dinhly}

\begin{dinhnghia}[title={Chuẩn toán tử}]
  \[
    \|T\| = \sup_{\|x\| \le 1} \|Tx\| = \sup_{x \ne 0} \frac{\|Tx\|}{\|x\|} = \inf\{C : \|Tx\| \le C\|x\|\}.
  \]
\end{dinhnghia}

\begin{dinhnghia}[title={Không gian $\mathcal{B}(X, Y)$}]
  Tập tất cả toán tử tuyến tính bị chặn $T: X \to Y$ với chuẩn toán tử. Nếu $Y$ là Banach thì $\mathcal{B}(X, Y)$ là Banach.
\end{dinhnghia}

\begin{dinhnghia}[title={Không gian đối ngẫu}]
  $X^* = \mathcal{B}(X, \mathbb{K})$ --- không gian các \textbf{phiếm hàm tuyến tính liên tục} (functional) trên $X$. $X^*$ luôn là Banach.
\end{dinhnghia}

%% ================================================================
%%  CHƯƠNG 4: CÁC ĐỊNH LÝ CƠ BẢN
%% ================================================================
\section{Các định lý cơ bản của giải tích hàm}

\begin{dinhly}[title={Định lý Hahn--Banach (dạng giải tích)}]
  Cho $X$ là không gian định chuẩn, $M \subset X$ là không gian con, $f_0 \in M^*$. Thì tồn tại $f \in X^*$ sao cho:
  \begin{enumerate}
    \item $f|_M = f_0$ (mở rộng).
    \item $\|f\| = \|f_0\|$ (bảo toàn chuẩn).
  \end{enumerate}
\end{dinhly}

\begin{tinhchat}[title={Hệ quả Hahn--Banach}]
  \begin{itemize}
    \item $\forall x \ne 0$, $\exists f \in X^*$: $\|f\| = 1$ và $f(x) = \|x\|$.
    \item $X^*$ phân li các điểm của $X$: nếu $f(x) = f(y)$ $\forall f \in X^*$ thì $x = y$.
    \item $\|x\| = \sup\{|f(x)| : f \in X^*,\, \|f\| \le 1\}$.
  \end{itemize}
\end{tinhchat}

\begin{dinhly}[title={Nguyên lý bị chặn đều (Banach--Steinhaus)}]
  Cho $X$ Banach, $Y$ định chuẩn, $(T_\alpha)_{\alpha \in A} \subset \mathcal{B}(X, Y)$. Nếu $\sup_\alpha \|T_\alpha x\| < \infty$ với mọi $x \in X$ thì $\sup_\alpha \|T_\alpha\| < \infty$.
\end{dinhly}

\begin{dinhly}[title={Định lý ánh xạ mở (Banach)}]
  Nếu $T: X \to Y$ là toán tử tuyến tính bị chặn toàn ánh giữa hai không gian Banach thì $T$ là ánh xạ mở (biến tập mở thành tập mở).

  Hệ quả: nếu $T$ song ánh thì $T^{-1}$ bị chặn (liên tục).
\end{dinhly}

\begin{dinhly}[title={Định lý đồ thị đóng}]
  $T: X \to Y$ tuyến tính ($X, Y$ Banach). $T$ bị chặn khi và chỉ khi đồ thị $\text{graph}(T) = \{(x, Tx) : x \in X\}$ là tập đóng trong $X \times Y$.
\end{dinhly}

%% ================================================================
%%  CHƯƠNG 5: KHÔNG GIAN HILBERT
%% ================================================================
\section{Không gian Hilbert}

\begin{dinhnghia}[title={Tích vô hướng}]
  \textbf{Không gian tích vô hướng}: không gian vectơ $H$ (trên $\mathbb{R}$ hoặc $\mathbb{C}$) với hàm $\langle \cdot, \cdot \rangle: H \times H \to \mathbb{K}$ thỏa:
  \begin{enumerate}
    \item $\langle x, x \rangle \ge 0$; \;$\langle x, x \rangle = 0 \Leftrightarrow x = 0$.
    \item $\langle \alpha x + \beta y, z \rangle = \alpha\langle x, z \rangle + \beta\langle y, z \rangle$ (tuyến tính vế trái).
    \item $\langle x, y \rangle = \overline{\langle y, x \rangle}$ (liên hợp đối xứng).
  \end{enumerate}
  Chuẩn: $\|x\| = \sqrt{\langle x, x \rangle}$.
\end{dinhnghia}

\begin{dinhnghia}[title={Không gian Hilbert}]
  Không gian tích vô hướng \textbf{đầy đủ} gọi là \textbf{không gian Hilbert}.
\end{dinhnghia}

\begin{dinhly}[title={Bất đẳng thức Cauchy--Schwarz}]
  $|\langle x, y \rangle| \le \|x\|\,\|y\|$. Đẳng thức xảy ra $\Leftrightarrow$ $x, y$ phụ thuộc tuyến tính.
\end{dinhly}

\begin{dinhnghia}[title={Trực giao}]
  $x \perp y$ nếu $\langle x, y \rangle = 0$. Tập $A \subset H$: $A^\perp = \{x \in H : \langle x, a \rangle = 0,\, \forall a \in A\}$.
\end{dinhnghia}

\begin{dinhly}[title={Định lý hình chiếu}]
  Cho $M$ là không gian con \textbf{đóng} của không gian Hilbert $H$. Thì:
  \[
    H = M \oplus M^\perp \quad\text{(tổng trực tiếp trực giao)}.
  \]
  $\forall x \in H$, $\exists!$ $m \in M$: $\|x - m\| = \min_{y \in M}\|x - y\|$ (hình chiếu trực giao).
\end{dinhly}

\begin{dinhly}[title={Định lý biểu diễn Riesz}]
  Cho $H$ là Hilbert và $f \in H^*$. Thì $\exists!$ $y \in H$ sao cho $f(x) = \langle x, y \rangle$ $\forall x$ và $\|f\| = \|y\|$.

  Vậy $H^* \cong H$ (đẳng cự phản tuyến tính).
\end{dinhly}

\subsection{Hệ trực chuẩn}

\begin{dinhnghia}
  $\{e_\alpha\}$ là \textbf{hệ trực chuẩn} nếu $\langle e_\alpha, e_\beta \rangle = \delta_{\alpha\beta}$. Hệ trực chuẩn là \textbf{đầy đủ} (cơ sở Hilbert) nếu $\overline{\text{span}\{e_\alpha\}} = H$.
\end{dinhnghia}

\begin{dinhly}[title={Khai triển Fourier tổng quát}]
  Cho $\{e_n\}$ là cơ sở trực chuẩn đếm được. Với $x \in H$:
  \[
    x = \sum_{n=1}^{\infty} \langle x, e_n \rangle\, e_n, \qquad \|x\|^2 = \sum_{n=1}^{\infty} |\langle x, e_n \rangle|^2 \;\text{(Parseval)}.
  \]
\end{dinhly}

\begin{dinhly}[title={Bất đẳng thức Bessel}]
  Với hệ trực chuẩn $\{e_n\}$ (không cần đầy đủ): $\displaystyle\sum_{n=1}^{\infty} |\langle x, e_n \rangle|^2 \le \|x\|^2$.

  Đẳng thức (Parseval) $\Leftrightarrow$ hệ đầy đủ.
\end{dinhly}

\begin{phuongphap}[title={Trực chuẩn hóa Gram--Schmidt}]
  Cho $\{v_1, v_2, \ldots\}$ độc lập tuyến tính:
  \[
    e_1 = \frac{v_1}{\|v_1\|}, \qquad e_n = \frac{v_n - \sum_{k=1}^{n-1}\langle v_n, e_k \rangle e_k}{\left\|v_n - \sum_{k=1}^{n-1}\langle v_n, e_k \rangle e_k\right\|}.
  \]
\end{phuongphap}

%% ================================================================
%%  CHƯƠNG 6: TOÁN TỬ TRÊN KHÔNG GIAN HILBERT
%% ================================================================
\section{Toán tử trên không gian Hilbert}

\begin{dinhnghia}[title={Toán tử liên hợp}]
  Cho $T \in \mathcal{B}(H)$. \textbf{Toán tử liên hợp} $T^*$ xác định bởi:
  \[
    \langle Tx, y \rangle = \langle x, T^*y \rangle, \qquad \forall x, y \in H.
  \]
  $T^*$ tồn tại duy nhất (Riesz) và $\|T^*\| = \|T\|$.
\end{dinhnghia}

\begin{dinhnghia}[title={Phân loại toán tử}]
  \begin{itemize}
    \item \textbf{Tự liên hợp} (self-adjoint): $T^* = T$ $\Leftrightarrow$ $\langle Tx, y \rangle = \langle x, Ty \rangle$.
    \item \textbf{Chuẩn tắc} (normal): $T^*T = TT^*$.
    \item \textbf{Unitary:} $T^*T = TT^* = I$ $\Leftrightarrow$ $T$ là đẳng cự toàn ánh.
    \item \textbf{Dương} (positive): $\langle Tx, x \rangle \ge 0$ $\forall x$.
    \item \textbf{Toán tử chiếu trực giao}: $P^2 = P = P^*$.
  \end{itemize}
\end{dinhnghia}

\begin{tinhchat}[title={Toán tử tự liên hợp}]
  Cho $T = T^* \in \mathcal{B}(H)$:
  \begin{itemize}
    \item Mọi trị riêng đều thực.
    \item Vectơ riêng ứng với trị riêng khác nhau thì trực giao.
    \item $\|T\| = \sup_{\|x\|=1} |\langle Tx, x \rangle|$.
  \end{itemize}
\end{tinhchat}

\subsection{Toán tử compact}

\begin{dinhnghia}
  $T: X \to Y$ là \textbf{compact} nếu $T$ biến mọi tập bị chặn thành tập tương đối compact (có bao đóng compact). Tương đương: mọi dãy bị chặn $(x_n)$, dãy $(Tx_n)$ có dãy con hội tụ.
\end{dinhnghia}

\begin{tinhchat}
  \begin{itemize}
    \item Toán tử compact là bị chặn (liên tục).
    \item Hợp toán tử compact với toán tử bị chặn là compact.
    \item Giới hạn (theo chuẩn toán tử) của dãy toán tử compact là compact.
    \item Toán tử hạng hữu hạn (ảnh hữu hạn chiều) là compact.
    \item Trên không gian vô hạn chiều, toán tử đồng nhất $I$ không compact.
  \end{itemize}
\end{tinhchat}

\begin{dinhly}[title={Lý thuyết Fredholm--Riesz}]
  Cho $T$ compact trên Banach $X$, $\lambda \ne 0$:
  \begin{enumerate}
    \item $\ker(\lambda I - T)$ hữu hạn chiều.
    \item $\text{Im}(\lambda I - T)$ đóng.
    \item $\lambda I - T$ toàn ánh $\Leftrightarrow$ $\lambda I - T$ đơn ánh.
    \item Phổ $\sigma(T) \setminus \{0\}$ là tập đếm được (hữu hạn hoặc dãy $\to 0$), mỗi $\lambda \ne 0$ trong phổ là trị riêng.
  \end{enumerate}
\end{dinhly}

\begin{dinhly}[title={Định lý phổ (Hilbert--Schmidt)}]
  Cho $T = T^*$ compact trên Hilbert $H$. Thì tồn tại cơ sở trực chuẩn $\{e_n\}$ gồm các vectơ riêng của $T$:
  \[
    Tx = \sum_{n=1}^{\infty} \lambda_n \langle x, e_n \rangle\, e_n,
  \]
  với $(\lambda_n)$ là dãy trị riêng thực, $\lambda_n \to 0$.
\end{dinhly}

%% ================================================================
%%  CHƯƠNG 7: LÝ THUYẾT PHỔ
%% ================================================================
\section{Lý thuyết phổ}

\begin{dinhnghia}
  Cho $T \in \mathcal{B}(X)$, $X$ Banach:
  \begin{itemize}
    \item \textbf{Tập phân giải} (resolvent set): $\rho(T) = \{\lambda \in \mathbb{C} : (\lambda I - T)^{-1} \in \mathcal{B}(X)\}$.
    \item \textbf{Phổ}: $\sigma(T) = \mathbb{C} \setminus \rho(T)$.
    \item \textbf{Toán tử phân giải}: $R_\lambda(T) = (\lambda I - T)^{-1}$, $\lambda \in \rho(T)$.
  \end{itemize}
\end{dinhnghia}

\begin{dinhnghia}[title={Phân loại phổ}]
  \begin{itemize}
    \item \textbf{Phổ điểm} $\sigma_p(T)$: $\lambda$ là trị riêng ($\ker(\lambda I - T) \ne \{0\}$).
    \item \textbf{Phổ liên tục} $\sigma_c(T)$: $(\lambda I - T)$ đơn ánh, $\overline{\text{Im}(\lambda I - T)} = X$ nhưng $\text{Im} \ne X$.
    \item \textbf{Phổ thặng dư} $\sigma_r(T)$: $(\lambda I - T)$ đơn ánh nhưng $\overline{\text{Im}(\lambda I - T)} \ne X$.
  \end{itemize}
\end{dinhnghia}

\begin{tinhchat}
  \begin{itemize}
    \item $\sigma(T)$ là tập compact, không rỗng.
    \item $\sigma(T) \subset \{\lambda : |\lambda| \le \|T\|\}$.
    \item \textbf{Bán kính phổ:} $r(T) = \max\{|\lambda| : \lambda \in \sigma(T)\} = \displaystyle\lim_{n \to \infty} \|T^n\|^{1/n}$.
    \item Toán tử phân giải $R_\lambda$ giải tích trên $\rho(T)$: $R_\lambda = \sum_{n=0}^{\infty} \dfrac{T^n}{\lambda^{n+1}}$ khi $|\lambda| > r(T)$.
  \end{itemize}
\end{tinhchat}

\begin{dinhly}[title={Phổ toán tử tự liên hợp}]
  Cho $T = T^*$ trên Hilbert:
  \begin{itemize}
    \item $\sigma(T) \subset \mathbb{R}$.
    \item $\sigma_r(T) = \varnothing$.
    \item $\sigma(T) \subset [m, M]$ với $m = \inf_{\|x\|=1}\langle Tx, x \rangle$, $M = \sup_{\|x\|=1}\langle Tx, x \rangle$.
    \item $m, M \in \sigma(T)$.
  \end{itemize}
\end{dinhly}

\begin{dinhly}[title={Phổ toán tử chuẩn tắc}]
  Cho $T$ chuẩn tắc ($T^*T = TT^*$) trên Hilbert: $\sigma_r(T) = \varnothing$ và $\|T\| = r(T)$.
\end{dinhly}

%% ================================================================
%%  CHƯƠNG 8: CÁC TÔPÔ YẾU
%% ================================================================
\section{Các tôpô yếu}

\begin{dinhnghia}
  \begin{itemize}
    \item \textbf{Tôpô yếu} trên $X$: tôpô yếu nhất sao cho mọi $f \in X^*$ liên tục. $x_n \xrightarrow{w} x$ nếu $f(x_n) \to f(x)$ $\forall f \in X^*$.
    \item \textbf{Tôpô yếu-$*$} trên $X^*$: $f_n \xrightarrow{w^*} f$ nếu $f_n(x) \to f(x)$ $\forall x \in X$.
  \end{itemize}
\end{dinhnghia}

\begin{dinhly}[title={Định lý Banach--Alaoglu}]
  Hình cầu đơn vị đóng $\{f \in X^* : \|f\| \le 1\}$ là compact trong tôpô yếu-$*$.
\end{dinhly}

\begin{dinhly}[title={Không gian phản xạ}]
  $X$ \textbf{phản xạ} (reflexive) nếu ánh xạ chính tắc $J: X \to X^{**}$, $Jx(f) = f(x)$ là đẳng cự toàn ánh.
  \begin{itemize}
    \item Hilbert $\Rightarrow$ phản xạ. $L^p$ ($1 < p < \infty$) phản xạ.
    \item $L^1$, $L^\infty$, $C[a,b]$ không phản xạ.
  \end{itemize}
\end{dinhly}

\begin{dinhly}[title={Eberlein--Šmulian}]
  Trong Banach phản xạ, hình cầu đơn vị đóng compact yếu (tương đương: mọi dãy bị chặn có dãy con hội tụ yếu).
\end{dinhly}

\end{document}